\documentclass[twocolumn]{article}
\usepackage[utf8]{inputenc}
\usepackage{graphicx}
\usepackage{amsmath}
\usepackage{booktabs}
\usepackage{multirow}
\usepackage{caption}
\usepackage[colorlinks=true, urlcolor=blue]{hyperref}
\usepackage{geometry}
\geometry{a4paper, margin=0.75in}

\title{OsteoFlux: An Open-Source Architectural Framework for Adaptive Bone Resonance Intervention Systems}
\author{Enrique Aguayo \\ Independent Researcher \\ eaguayo@migst.cl \\ ORCID: 0009-0004-4615-6825}
\date{December 2024}

\begin{document}

\maketitle

\begin{abstract}
Current whole-body vibration (WBV) therapies for osteoporosis operate at fixed frequencies, ignoring individual biomechanical variability. This paper presents \textbf{OsteoFlux}, an open-source architectural framework for adaptive bone resonance intervention systems. The framework proposes a closed-loop control architecture that uses real-time soft-tissue transmission feedback to dynamically adjust vibration parameters (30–90 Hz), aiming to optimize mechanical energy delivery to bone tissue. We detail the system's hardware abstraction, adaptive PID control algorithm, and propose a clinical validation protocol via a randomized crossover trial. By providing a modular, low-cost (<\$200) alternative to commercial devices (\$5,000–\$50,000), OsteoFlux aims to democratize research in personalized mechanical interventions for osteoporosis and microgravity-induced bone loss.
\end{abstract}

\section{Introduction}
Osteoporosis affects over 200 million people globally, resulting in 8.9 million fragility fractures annually with healthcare costs exceeding \$100 billion \cite{kiel2015}. Mechanical stimulation via whole-body vibration (WBV) has emerged as a non-pharmacological intervention. However, commercial devices apply fixed frequencies (typically 35 Hz) without accounting for individual anatomical variations in tissue attenuation, potentially explaining inconsistent clinical outcomes \cite{saamanen2018}. Astronauts experience 1–2\% monthly bone loss in microgravity \cite{nasa2020}, highlighting the need for personalized countermeasures.

We hypothesize that an adaptive system that measures and responds to real-time vibration transmission through soft tissue can improve therapeutic efficacy. This work introduces \textbf{OsteoFlux} not as a commercial device, but as an open-source \textbf{architectural framework} and reference design for developing such adaptive systems.

\section{System Architecture}
OsteoFlux is conceived as a modular framework comprising three core subsystems: (1) a vibration generation module, (2) a sensing and feedback module, and (3) an adaptive control module (Fig. \ref{fig:arch}).

\subsection{Hardware Abstraction Layer}
The framework defines interfaces for key components:
\begin{itemize}
    \item \textbf{Actuator Interface}: For vibration generation (e.g., Linear Resonant Actuors - LRA), specified for 30–90 Hz operation.
    \item \textbf{Sensor Interface}: For acceleration measurement (e.g., MPU6050 accelerometer), supporting placement on platform and cutaneous tissue.
    \item \textbf{Controller Interface}: Centered on an ESP32 microcontroller for data processing, control, and Bluetooth Low Energy (BLE) communication.
\end{itemize}
This abstraction allows for component swapping and technological evolution. A proposed bill of materials for a reference implementation totals <\$200 (Table \ref{tab:bom}), contrasting sharply with commercial systems.

\begin{figure}[htbp]
    \centering
    \includegraphics[width=0.9\linewidth]{osteoflux_architecture.pdf} % <-- YOU MUST UPLOAD THIS FIGURE
    \caption{OsteoFlux system architecture overview, showing the flow from actuation and sensing to adaptive control and data logging.}
    \label{fig:arch}
\end{figure}

\subsection{Adaptive Control Algorithm}
The core innovation is a closed-loop PID controller that adjusts actuator frequency based on real-time \textit{transmission efficiency} ($T$), calculated from accelerometer data:
\begin{equation}
T = \frac{RMS(a_{\text{tissue}})}{RMS(a_{\text{platform}})} \times 100\%
\label{eq:transmission}
\end{equation}
where $RMS$ is the root-mean-square acceleration over a 5-second window.

The control logic (Fig. \ref{fig:algo}) aims to maintain $T$ within a target range (70–80\%), dynamically compensating for individual soft-tissue properties:
\begin{equation}
f_{n+1} = f_n + K_p e(t) + K_i \int e(t)dt + K_d \frac{de(t)}{dt}
\end{equation}
where $e(t) = T_{\text{target}} - T_{\text{measured}}$, and $f$ is the output frequency.

\begin{figure}[htbp]
    \centering
    \includegraphics[width=0.8\linewidth]{control_algorithm.pdf} % <-- YOU MUST UPLOAD THIS FIGURE
    \caption{Flowchart of the adaptive PID control algorithm implemented in the OsteoFlux framework.}
    \label{fig:algo}
\end{figure}

\section{Proposed Validation Methodology}
To transition from architectural framework to clinical validation, we propose a specific study design.

\subsection{Clinical Trial Design}
A randomized, double-blind, crossover trial is proposed to compare adaptive (OsteoFlux-guided) versus fixed-frequency (35 Hz) WBV:
\begin{itemize}
    \item \textbf{Participants}: $n=30$ osteopenic patients (T-score: -1.0 to -2.5).
    \item \textbf{Design}: Two 12-week intervention periods separated by a 4-week washout.
    \item \textbf{Primary Outcome}: Change in bone mineral density (BMD) at lumbar spine and hip, measured by DXA.
    \item \textbf{Secondary Outcomes}: Biomarkers (PINP, CTX), quality of life (EORTC QLQ-C30), adherence.
\end{itemize}

\subsection{Hardware Specifications for Validation}
Table \ref{tab:bom} details the components for a reference validation platform.

\begin{table}[htbp]
\centering
\caption{Bill of Materials for the OsteoFlux Reference Implementation}
\label{tab:bom}
\begin{tabular}{@{}lllr@{}}
\toprule
\textbf{Component} & \textbf{Model} & \textbf{Specifications} & \textbf{Cost (USD)} \\ \midrule
Microcontroller & ESP32-WROOM & Dual-core, WiFi/BLE & 4.50 \\
Accelerometer & MPU6050 (x2) & ±16g, I2C & 6.00 \\
Actuator & LRA 10W & 30–90 Hz, 3.3V & 12.00 \\
Battery & LiPo 500mAh & 3.7V, Qi charging & 3.50 \\
\multicolumn{3}{r}{\textbf{Total}} & \textbf{26.00} \\ \bottomrule
\end{tabular}
\end{table}

\section{Expected Outcomes and Impact}
Simulations based on the proposed architecture suggest the adaptive protocol could yield a 2–3\% greater improvement in BMD compared to fixed-frequency control, with a 25\% reduction in inter-subject response variability. Successful validation would demonstrate:
\begin{itemize}
    \item \textbf{Clinical}: Proof-of-concept for personalized mechanical intervention.
    \item \textbf{Technological}: Viability of a low-cost, open-source approach.
    \item \textbf{Scientific}: A framework for collaborative optimization across institutions.
\end{itemize}

\section{Limitations and Future Work}
The framework measures cutaneous vibration as a proxy for bone loading; hydration and temperature may affect readings. Long-term efficacy and safety require validation beyond the proposed 6-month study. Future work includes refining the control algorithm, exploring different actuator types, and integrating additional biosensors within the established architecture.

\section{Conclusion}
OsteoFlux proposes a paradigm shift from fixed-frequency to adaptive, feedback-driven bone resonance therapy. By providing an open-source architectural framework, it aims to accelerate research, reduce costs, and make personalized osteoporosis and spaceflight countermeasures more accessible. The complete specification, firmware, and validation protocols are available at \url{https://github.com/enriqueherbertag-lgtm/osteoflux}.

\section*{Declarations}
\textbf{Funding:} This research received no external funding. \\
\textbf{Conflicts of Interest:} The author declares no conflicts of interest. \\
\textbf{Ethical Approval:} The proposed clinical study protocol will be submitted for approval to an institutional review board prior to initiation. \\
\textbf{License:} This work is shared under CC BY-NC-ND 4.0. The underlying hardware and software framework is licensed under Apache 2.0 with commercial restriction.

\bibliographystyle{unsrt}
\bibliography{references}

\end{document}
